\section*{第 \ref{ch:b1:quantum-introduction} 章——量子物理导论}

\begin{solution}{exo:b1:quantum-introduction:1}
$h\nu$：\qty{100}{MHz}：$6.6 \times 10^{-26}$~J $= \qty{4.1e-7}{eV}$；\qty{500}{nm}：
$1240/500 = \qty{2.5}{eV}$；\qty{0.10}{nm}：\qty{12.4}{keV}；\qty{1}{MeV}：$\lambda =
1240/10^6 = \qty{1.2e-3}{nm} = \qty{1.2}{pm}$。激光：$1240/633 = \qty{1.96}{eV} =
\qty{3.1e-19}{J}$，因此 $10^{-3}/3.1 \times 10^{-19} = 3.2 \times 10^{15}$ 个光子
每秒。
\end{solution}

\begin{solution}{exo:b1:quantum-introduction:2}
$\lambda_0 = 1240/2.3 = \qty{540}{nm}$。在 \qty{400}{nm} 处，$h\nu = \qty{3.1}{eV}$：
$E_{k,\max} = \qty{0.8}{eV}$，$V_s = \qty{0.8}{V}$。光强加倍，
电流加倍，$V_s$ 保持不变。在 \qty{600}{nm} 处
（\qty{2.07}{eV} $< W$）：无论光强多大，都没有光电子发射。
\end{solution}

\begin{solution}{exo:b1:quantum-introduction:3}
电子：$1.23/\sqrt{100} = \qty{0.12}{nm}$；在 \qty{1}{keV} 时：\qty{0.039}{nm}。
中子：$E = \qty{4.0e-21}{J}$，$p = \sqrt{2m_nE} = \qty{3.7e-24}{kg.m/s}$，
$\lambda = h/p = \qty{0.18}{nm}$。小球：$6.6 \times 10^{-34}/10^{-3} = \qty{7e-31}{m}$。
前三者的 $\lambda$ 与原子间距同数量级，
均会发生衍射（电子、中子及 X 射线晶体学）；小球则永不发生。
\end{solution}

\begin{solution}{exo:b1:quantum-introduction:4}
$E_1 = h^2/8m_eL^2 = 4.39 \times 10^{-67}/(8 \times 9.11 \times 10^{-31} \times 10^{-18}) =
\qty{6.0e-20}{J} = \qty{0.38}{eV}$，$E_2 = 4E_1 = \qty{1.5}{eV}$；光子
\qty{1.13}{eV}，$\lambda = 1240/1.13 = \qty{1.1}{\micro m}$（近红外）。质子
限制在 \qty{5}{fm} 内：$4.39 \times 10^{-67}/(8 \times 1.67 \times 10^{-27} \times 25 \times 10^{-30}) =
\qty{1.3e-12}{J} = \qty{8}{MeV}$ ——核能的量级，
仅由束缚即可得到。
\end{solution}

\begin{solution}{exo:b1:quantum-introduction:5}
斜率 $(2.05 - 0.40)/(4.0 \times 10^{14}) = \qty{4.1e-15}{V.s}$，$h = e \times 4.1 \times
10^{-15} = \qty{6.6e-34}{J.s}$。$W = h\nu - eV_s = 2.48 - 0.40 = \qty{2.1}{eV}$，
截止波长 $\lambda_0 = 1240/2.1 = \qty{600}{nm}$：金属为铯。
\end{solution}

\begin{solution}{exo:b1:quantum-introduction:6}
$\Delta\lambda = \qty{2.43}{pm}$：$\lambda' = \qty{12.4}{pm}$。散射前：$1240/0.010 =
\qty{124}{keV}$；散射后：$1240/0.0124 = \qty{100}{keV}$；电子获得
\qty{24}{keV}。对于光，$\Delta\lambda/\lambda = 2.4 \times 10^{-12}/5 \times 10^{-7} = 5 \times
10^{-6}$：1923 年无法测量——X 射线使该偏移成为 $25\%$ 的效应。
\end{solution}

\begin{solution}{exo:b1:quantum-introduction:7}
$p = \sqrt{2 \times 9.11 \times 10^{-31} \times 5 \times 10^4 \times 1.6 \times 10^{-19}} =
\qty{1.2e-22}{kg.m/s}$，$\lambda = \qty{5.5}{pm}$（相对论修正后为
\qty{5.4}{pm}）；条纹间距 $\lambda D/d = 5.5 \times 10^{-12} \times 0.35/2 \times 10^{-6}
= \qty{1}{\micro m}$：肉眼不可见，因此需要放大透镜。十个
电子：十个点，看似随机分布；$10^5$ 个：出现条纹。每个
电子只落在一点上；落点的\emph{概率}遵循
双缝波动图样——单个电子的波同时通过两条缝。
\end{solution}

\begin{solution}{exo:b1:quantum-introduction:8}
原子：$\Delta p \geq \hbar/2\Delta x = \qty{5.3e-25}{kg.m/s}$，$E \sim \Delta p^2/2m_e =
\qty{1.5e-19}{J} \approx \qty{1}{eV}$：量级正确。原子核：$\Delta p \geq
\qty{1.1e-20}{kg.m/s}$；$\Delta p^2/2m_e$ 将达到 \qty{400}{MeV}，远高于
$m_ec^2$，因此取 $E \approx \Delta p\,c = \qty{3e-12}{J} = \qty{20}{MeV}$：原子核内任何力
都无法束缚如此高能的电子——原子核
中不含电子（β 电子是在发射时产生的）。颗粒：$\Delta v \geq
\hbar/2m\Delta x = 1.05 \times 10^{-34}/(2 \times 10^{-9} \times 10^{-6}) = \qty{5e-20}{m/s}$：
可忽略。
\end{solution}

\begin{solution}{exo:b1:quantum-introduction:9}
(a) $m_ev^2/r = e^2/4\pi\varepsilon_0r^2$：$v = \sqrt{e^2/4\pi\varepsilon_0m_er}$，$E =
\tfrac12m_ev^2 - e^2/4\pi\varepsilon_0r = -e^2/8\pi\varepsilon_0r$。(b) $m_evr = n\hbar$ 且
$m_ev^2r = e^2/4\pi\varepsilon_0$：将第一式平方后除以第二式：
$m_er = n^2\hbar^2 \cdot 4\pi\varepsilon_0/e^2$，$r_n = n^2a_0$，且 $E_n = -e^2/8\pi\varepsilon_0n^2a_0
= E_1/n^2$，其中 $a_0 = \qty{52.9}{pm}$，$E_1 = \qty{-13.6}{eV}$。(c) $v_1 =
\hbar/m_ea_0 = \qty{2.19e6}{m/s}$，$v_1/c = 1/137$。(d) $2 \to 1$：$\tfrac34 \times
13.6 = \qty{10.2}{eV}$，\qty{122}{nm}；$3 \to 2$：$(\tfrac14 - \tfrac19) \times 13.6 =
\qty{1.89}{eV}$，\qty{656}{nm}。(e) $\lambda_n = h/m_ev_n = nh/m_ev_1 = 2\pi na_0$
（因为 $m_ev_1a_0 = \hbar$），且 $2\pi r_n = 2\pi n^2a_0 = n\lambda_n$。
\end{solution}

\begin{solution}{exo:b1:quantum-introduction:10}
(a) 将 $e^2$ 替换为 $Ze^2$：$a_0 \to a_0/Z$，$E_1 \to Z^2E_1$。He$^+$：
\qty{54.4}{eV}。U$^{91+}$：$92^2 \times 13.6 = \qty{115}{keV}$——但此时 $v_1 = Z\alpha c
= 0.67c$：此处非相对论模型仅具指示意义。(b) $a_0/207
= \qty{256}{fm}$，$E_1 = 207 \times 13.6 = \qty{2.8}{keV}$；赖曼 $\alpha$ 线：
\qty{2.1}{keV}，$\lambda = \qty{0.59}{nm}$（X 射线）。μ 子的轨道仅约为
$300$ 个质子半径：其能级随质子大小可测地移动，
这正是（2010 年）重新测定质子半径的方法。(c) 两个
粒子绕其共同质心运动：约化质量 $m_e/2$ 使
每一能量减半：$E_1 = \qty{-6.8}{eV}$，赖曼 $\alpha$ 线位于 \qty{5.1}{eV}，\qty{243}{nm}。
\end{solution}

\begin{solution}{exo:b1:quantum-introduction:11}
(a) $h^2/8m_e = \qty{6.02e-38}{J.m^2}$。$R = \qty{2}{nm}$：$E_e = 6.02 \times 10^{-38}/
(0.13 \times 4 \times 10^{-18}) = \qty{0.72}{eV}$，$E_h = \qty{0.21}{eV}$，总能量 $1.74 +
0.93 = \qty{2.67}{eV}$，\qty{464}{nm}：蓝色。\qty{3}{nm}：束缚能 $\times4/9$，
\qty{2.15}{eV}，\qty{577}{nm}：黄色。\qty{4}{nm}：$\times1/4$，\qty{1.97}{eV}，
\qty{629}{nm}：红色。(b) \qty{520}{nm} 对应 \qty{2.38}{eV}：束缚能 \qty{0.65}{eV}
$= 0.93 \times (2/R)^2$，$R = \qty{2.4}{nm}$。(c) 无束缚时：仅 $E_g$，
\qty{713}{nm}。(d) $\delta E = -2 \times 0.1 \times 0.41 = \qty{-0.08}{eV}$，占
\qty{2.15}{eV} 的 $4\%$：约 \qty{20}{nm} 的波长展宽——纯色需要
单分散的样品批次。
\end{solution}

\begin{solution}{exo:b1:quantum-introduction:12}
(a) $E(\Delta x) = \hbar^2/8m\Delta x^2 + \tfrac12m\omega^2\Delta x^2$；导数为零时
$\Delta x^2 = \hbar/2m\omega$，此时两项均等于 $\hbar\omega/4$：$E_{\min} =
\hbar\omega/2$。(b) $\tfrac12 \times 1.055 \times 10^{-34} \times 8.3 \times 10^{14} = \qty{4.4e-20}{J}
= \qty{0.27}{eV}$；$T \sim \hbar\omega/k_B = \qty{6300}{K}$——在室温下
振动冻结于其基态。(c) $\Delta p = \hbar/2\Delta x =
\qty{1.05e-24}{kg.m/s}$，$E = \Delta p^2/2m_{\mathrm{He}} = 1.1 \times 10^{-48}/1.33 \times
10^{-26} = \qty{8e-23}{J} = \qty{0.5}{meV}$：约为束缚能的一半——原子无法
静止到足以结晶的程度；氦在
绝对零度附近仍保持液态，除非受压（\qty{25}{bar}）。(d) $E_1 = \pi^2\hbar^2/2mL^2$ 与
$\hbar^2/2mL^2$ 相比：大十倍，理应如此——该束缚只是
下界，而盒中态的 $\Delta x \approx 0.18L$，而非 $L/2$。
\end{solution}

\begin{solution}{pb:b1:quantum-introduction:1}
\textbf{1.} 光使阴极逸出能量各异的电子；阳极保持在 $-V_s$，
于是阳极排斥电子，只有
$E_k > eV_s$ 的电子能够到达。电流为零的条件是 $eV_s = E_{k,\max}$。

\textbf{2.} $eV_s = h\nu - W$：斜率为 $h/e$，截距为 $-W/e$；$V_s = 0$ 时
即截止频率 $\nu_0 = W/h$。

\textbf{3.} 取两端点：$(2.36 - 0.30)/(5.0 \times 10^{14}) = \qty{4.12e-15}{V.s}$，
$h = 1.602 \times 10^{-19} \times 4.12 \times 10^{-15} = \qty{6.60e-34}{J.s}$（对全部五个点拟合得
$6.62$）；误差在 $0.5\%$ 以内（参照值 $6.626 \times 10^{-34}$）。

\textbf{4.} $W = h\nu - eV_s = 2.47 - 0.30 = \qty{2.2}{eV}$；$\nu_0 = W/h =
\qty{5.3e14}{Hz}$，$\lambda_0 = 1240/2.2 = \qty{560}{nm}$。红色激光笔
（\qty{650}{nm}，\qty{1.9}{eV}）无法使电子逸出。

\textbf{5.} $V_s$ 不变而电流加倍：光子数加倍，每个光子的
能量不变。若为波动，则光越强，每个电子获得的能量越大，
且阈值应出现在强度而非频率上。

\textbf{6.} 原子所接收的功率为 $10^{-6} \times 10^{-20} = \qty{1e-26}{W}$；\qty{2.2}{eV} $=
\qty{3.5e-19}{J}$ 需要 $3.5 \times 10^7$~s——约一年。观测到的发射
在纳秒之内即可发生：能量以完整的量子一份份到达。

\textbf{7.} $m_ev^2/r = e^2/4\pi\varepsilon_0r^2$：$v = \sqrt{e^2/4\pi\varepsilon_0m_er}$；$E =
\tfrac12m_ev^2 - e^2/4\pi\varepsilon_0r = -e^2/8\pi\varepsilon_0r$。

\textbf{8.} 由 $(m_evr)^2 = n^2\hbar^2$ 与 $m_ev^2r = e^2/4\pi\varepsilon_0$ 联立：$r_n = n^2 \cdot
4\pi\varepsilon_0\hbar^2/m_ee^2 = n^2a_0$，$a_0 = (1.055 \times 10^{-34})^2/(8.99 \times 10^9
\times 9.11 \times 10^{-31} \times 2.57 \times 10^{-38}) = \qty{52.9}{pm}$。

\textbf{9.} $E_n = -e^2/8\pi\varepsilon_0r_n = -m_ee^4/2(4\pi\varepsilon_0)^2\hbar^2n^2$；$E_1 =
-2.31 \times 10^{-28}/(2 \times 5.29 \times 10^{-11}) = \qty{-2.18e-18}{J} = \qty{-13.6}{eV}$。

\textbf{10.} $v_1 = \hbar/m_ea_0 = \qty{2.19e6}{m/s}$；$v_1/c = 7.3 \times 10^{-3} =
1/137$。

\textbf{11.} $2 \to 1$：\qty{10.2}{eV}，\qty{122}{nm}（紫外线）；$3 \to 2$：
\qty{1.89}{eV}，\qty{656}{nm}（红光）；$4 \to 2$：\qty{2.55}{eV}，\qty{486}{nm}
（蓝绿色）：可见的正是巴耳末线。

\textbf{12.} $hc/\lambda = |E_1|(1/n^2 - 1/p^2)$：$R_H = |E_1|/hc = 13.6/1240 =
\qty{1.097e-2}{nm^{-1}} = \qty{1.097e7}{m^{-1}}$。巴耳末极限（$p \to \infty$）：
$\lambda = 4/R_H = \qty{365}{nm}$。

\textbf{13.} \qty{13.6}{eV}；$\lambda = 1240/13.6 = \qty{91}{nm}$。

\textbf{14.} $v_n = v_1/n$，$\lambda_n = h/m_ev_n = nh/m_ev_1 = 2\pi na_0$（因 $m_ev_1a_0
= \hbar$）；$2\pi r_n = 2\pi n^2a_0 = n\lambda_n$：轨道上恰好容纳 $n$ 个完整
波——玻尔的量子化规则即驻波规则。

\textbf{15.} 错误之处：电子并无轨道（它会辐射；基态的角动量为零），且该模型
对任何含两个电子的原子都失效。正确之处在于：氢原子的能级，进而
其光谱的每一条谱线。轨道被波函数取代，
即一团概率云，
其尺度为 $\sim n^2a_0$。

\textbf{16.} 驻波：$L = n\lambda/2$，$p_n = h/\lambda_n = nh/2L$，$E_n =
p_n^2/2m = n^2h^2/8mL^2$。球形情形只需将 $L \to R$ 代入：
基态在半径内恰好容纳半个波长。

\textbf{17.} $E_e = 6.02 \times 10^{-38}/(0.13 \times 4 \times 10^{-18}) = \qty{0.72}{eV}$；
$E_h = 0.13/0.45 \times 0.72 = \qty{0.21}{eV}$。

\textbf{18.} $R = \qty{2}{nm}$：$1.74 + 0.93 = \qty{2.67}{eV}$，\qty{464}{nm}，蓝色；
\qty{3}{nm}：\qty{2.15}{eV}，\qty{577}{nm}，黄色；\qty{4}{nm}：\qty{1.97}{eV}，
\qty{629}{nm}，红色。

\textbf{19.} 束缚能随 $1/R^2$ 增长：较小的盒子抬高
能级，光子能量更大，光也更偏蓝。

\textbf{20.} $\Delta p \sim \hbar/R = \qty{5.3e-26}{kg.m/s}$，$E \sim \Delta p^2/2m_e^\ast =
\qty{1.2e-20}{J} = \qty{0.07}{eV}$：量级正确（精确系数
$\pi^2/2$ 与较粗略的 $\Delta p$ 解释了十倍之差）。

\textbf{21.} $1240/577 = \qty{2.15}{eV} = \qty{3.4e-19}{J}$：$10^{-9}/3.4 \times 10^{-19}
= 2.9 \times 10^9$ 个光子每秒。

\textbf{22.} $10^{-3}/(1.96 \times 1.6 \times 10^{-19}) = 3.2 \times 10^{15}$ 每秒。

\textbf{23.} $100 \times 2.48 \times 1.6 \times 10^{-19} = \qty{4e-17}{J}$，在 \qty{0.1}{s} 内：
\qty{4e-16}{W}。

\textbf{24.} $h\nu = \qty{6.6e-26}{J}$：每秒 $1.5 \times 10^{13}$ 个光子——
每个光子都比天线的热能小 $10^5$ 倍（热能即 $k_BT$）：
探测不可能实现，也无必要；波动图像在此是精确的。

\textbf{25.} 原子的尺度 $a_0 = \qty{53}{pm}$；其能量标度，
\qty{13.6}{eV}（因而有电子伏特量级的光子、化学与可见光谱）；以及
电子的速率 $\alpha c \approx \qty{2e6}{m/s}$——这三者
均由 $h$、$e$ 与 $m_e$（及 $\varepsilon_0$）导出，没有任何人为调节。
\end{solution}
